% Source: https://artofproblemsolving.com/community/c3h1225424_2016_lmt_individual_14
% Stable identity: aops-topic-1225424-q14
In triangle $APB$, the point $Q$ on $AB$ divides it in the ratio $AP:PB$, so by the converse of the angle bisector theorem $PQ$ bisects $\angle APB$. A bisector that is also perpendicular to the opposite side makes the triangle isosceles: triangles $APQ$ and $BPQ$ are congruent (ASA), so $PA=PB$.

\begin{asy}
usepackage("amsmath");
size(3.5cm);
defaultpen(fontsize(10pt));
pair A=(0,0),C=(6,0),B=3*dir(70)+(3,0),P=(A+C)/2,Q=(A+B)/2;
draw(A--B--C--cycle,linewidth(1)); draw(P--B,blue); draw(P--Q,red);
label("$A$",A,SW);label("$B$",B,N);label("$C$",C,E);label("$P$",P,S);label("$Q$",Q,NW);
label("\mbox{}",(min(currentpicture,user=true)+max(currentpicture,user=true))/2,black); // keeps the text after the figure black when LaTeX typesets the labels inline
\end{asy}

Now $PA=PB=PC$, so $P$ is the circumcentre of $ABC$ and lies on side $AC$; thus $AC$ is a diameter and $\angle CBA=\boxed{90^\circ}$.
